\documentclass[11pt,a4paper]{article}

\usepackage[utf8]{inputenc}
\usepackage[T1]{fontenc}
\usepackage{amsmath,amssymb}
\usepackage{graphicx}
\usepackage{subcaption}
\usepackage{booktabs}
\usepackage{hyperref}
\usepackage{xcolor}
\usepackage{listings}
\usepackage{geometry}
\geometry{margin=2.5cm}

\lstset{
  basicstyle=\ttfamily\small,
  keywordstyle=\color{blue},
  commentstyle=\color{gray},
  breaklines=true,
  frame=single,
  language=C
}

\title{\textbf{Jacobi Iteration on a 2D Grid:\\
  Serial, OpenMP, and Hybrid MPI+OpenMP Parallelisation}}
\author{}
\date{March 2026}

\begin{document}

\maketitle
\tableofcontents
\newpage

% ============================================================
\section{The Numerical Problem}
% ============================================================

\subsection{Governing equation}

The benchmark solves the two-dimensional \emph{steady-state Laplace equation}
\begin{equation}
  \nabla^2 u(x,y)
  \;=\;
  \frac{\partial^2 u}{\partial x^2} + \frac{\partial^2 u}{\partial y^2}
  \;=\; 0,
  \qquad (x,y)\in\Omega \subset \mathbb{R}^2,
  \label{eq:laplace}
\end{equation}
subject to Dirichlet boundary conditions on the rectangular domain~$\Omega$.
Equation~\eqref{eq:laplace} models, among other phenomena, steady heat conduction
and electrostatic potential in source-free regions.

\subsection{Discretisation}

The domain is discretised on a uniform $N_x \times N_y$ grid with grid spacing
$h$ in both directions.  Including one layer of boundary (halo) cells, the
computational array has size $(N_x+2)\times(N_y+2)$ and is stored in row-major
(C) order as a flat 1D \texttt{double} array.  The linear index of cell $(i,j)$
is
\begin{equation}
  k(i,j) = i\,(N_y+2) + j,
  \qquad 0 \le i \le N_x+1,\quad 0 \le j \le N_y+1.
\end{equation}

\subsection{Jacobi iterative scheme}

At each iteration the interior values are updated via the five-point stencil
\begin{equation}
  u^{(t+1)}_{i,j}
  \;=\;
  \frac{1}{4}\!\left(
    u^{(t)}_{i-1,j} + u^{(t)}_{i+1,j}
    + u^{(t)}_{i,j-1} + u^{(t)}_{i,j+1}
  \right),
  \qquad 1 \le i \le N_x,\; 1 \le j \le N_y.
  \label{eq:jacobi}
\end{equation}
The method requires two grids (\texttt{grid} and \texttt{grid\_new}); after
each sweep the new values are copied back to the active grid.

\subsection{Boundary conditions}

The four sides are assigned fixed values:
\begin{center}
\begin{tabular}{ll}
\toprule
Boundary & Value \\
\midrule
Top ($i=0$)         & \texttt{TOP}    = 1.0 \\
Bottom ($i=N_x+1$)  & \texttt{BOTTOM} = 10.0 \\
Left ($j=0$)        & \texttt{LEFT}   = 1.0 \\
Right ($j=N_y+1$)   & \texttt{RIGHT}  = 1.0 \\
\bottomrule
\end{tabular}
\end{center}

\subsection{Convergence criterion}

Iterations stop when the relative residual norm falls below a tolerance
$\varepsilon = 10^{-4}$, or after a maximum of $5\times10^5$ iterations.
The residual is computed before each sweep as
\begin{equation}
  r^{(t)}
  \;=\;
  \sqrt{
    \sum_{i=1}^{N_x}\sum_{j=1}^{N_y}
    \left( 4\,u^{(t)}_{i,j}
           - u^{(t)}_{i-1,j} - u^{(t)}_{i+1,j}
           - u^{(t)}_{i,j-1} - u^{(t)}_{i,j+1}
    \right)^2
  },
  \label{eq:residual}
\end{equation}
and convergence is declared when $r^{(t)}/r^{(0)} < \varepsilon$.

% ============================================================
\section{Code Structure and Implementation}
% ============================================================

Three C99 executables are provided, sharing the same CLI
\texttt{./jacobi{\_}* <NX> <NY>}.
All functions are declared in a companion \texttt{.h} header so that the
\texttt{.c} translation unit contains only \texttt{main}.

\subsection{Serial baseline (\texttt{jacobi})}

The header \texttt{jacobi.h} provides four functions called in sequence:

\begin{enumerate}
  \item \textbf{\texttt{fill\_boundary\_grid}} --- sets the four boundary rows/columns
        in both \texttt{grid} and \texttt{grid\_new}.
  \item \textbf{\texttt{fill\_grid\_withouth\_boundary}} --- zero-initialises all
        interior cells.
  \item \textbf{\texttt{update}} --- performs one Jacobi sweep
        (equation~\eqref{eq:jacobi}).
  \item \textbf{\texttt{check\_convergence}} --- evaluates the sum of squared
        residuals (equation~\eqref{eq:residual}), without the square root, so
        that the caller can accumulate contributions from multiple sub-domains.
\end{enumerate}

Wall-clock timing is obtained via
\texttt{clock\_gettime(CLOCK\_MONOTONIC)}, which gives nanosecond resolution
and is unaffected by system-clock adjustments.

\subsection{OpenMP version (\texttt{jacobi\_just\_omp})}

The header \texttt{jacobi\_just\_omp.h} parallelises the two compute-intensive
kernels with OpenMP pragmas:

\begin{lstlisting}[language=C, caption={OpenMP loop parallelism in \texttt{update}.}]
#pragma omp parallel for private(i, j, k)
for (i = 1; i <= nx; i++) {
    for (j = 1; j <= ny; j++) {
        k = (ny+2)*i + j;
        grid_new[k] = 0.25 * (grid[k-1] + grid[k+1]
                             + grid[k-(ny+2)] + grid[k+(ny+2)]);
    }
}
\end{lstlisting}

\begin{lstlisting}[language=C, caption={Race-condition-safe reduction in \texttt{check\_convergence}.}]
#pragma omp parallel for private(i,j,k) reduction(+:tmpnorm)
for (i = 1; i <= nx; i++) {
    for (j = 1; j <= ny; j++) { ... tmpnorm += diff*diff; }
}
\end{lstlisting}

The \texttt{reduction(+:tmpnorm)} clause lets each thread accumulate into a
private copy; OpenMP merges the copies at the end of the parallel region,
eliminating any data race on the accumulator.
Timing is performed with \texttt{omp\_get\_wtime()}, called exclusively from the
master thread before and after each parallel region to avoid concurrent writes
to the timing variables.

\subsection{Hybrid MPI+OpenMP version (\texttt{jacobi\_omp\_mpi})}

\subsubsection{Domain decomposition}

The $N_x$ rows are partitioned into contiguous 1D strips among $P$ MPI ranks.
Rank $r$ owns
\begin{equation}
  \textit{local\_nx}(r)
  = \left\lfloor \tfrac{N_x}{P} \right\rfloor
  + \mathbf{1}[r < N_x \bmod P],
\end{equation}
so remainder rows are distributed among the lowest-rank processes, keeping
load imbalance to at most one row.
Each rank allocates
$(\textit{local\_nx}+2)\times(N_y+2)$ doubles, where the extra two rows are
\emph{ghost (halo) rows} storing data received from neighbouring ranks.

\subsubsection{Halo exchange}

After the convergence check and before each \texttt{update}, a non-blocking
halo exchange is performed:

\begin{lstlisting}[language=C, caption={Non-blocking halo exchange (excerpt).}]
MPI_Isend(&grid[1*(ny+2)],          ny+2, MPI_DOUBLE, rank-1, 0, ...);
MPI_Irecv(&grid[0*(ny+2)],          ny+2, MPI_DOUBLE, rank-1, 1, ...);
MPI_Isend(&grid[local_nx*(ny+2)],   ny+2, MPI_DOUBLE, rank+1, 1, ...);
MPI_Irecv(&grid[(local_nx+1)*(ny+2)], ny+2, MPI_DOUBLE, rank+1, 0, ...);
MPI_Waitall(req_cnt, reqs, MPI_STATUSES_IGNORE);
\end{lstlisting}

Using \texttt{MPI\_Isend}/\texttt{MPI\_Irecv} instead of blocking variants
avoids deadlock regardless of process ordering and allows the MPI library to
pipeline transfers.

\subsubsection{Global residual}

Each rank computes a \emph{local} squared residual (the sum of squared
differences over its own rows, without the square root).
\texttt{MPI\_Allreduce(MPI\_SUM)} accumulates these across all ranks;
the global square root is then taken exactly once:
\begin{equation}
  r^{(t)}_{\text{global}}
  = \sqrt{\sum_{p=0}^{P-1} r^{(t)}_{\text{local},\,p}}.
\end{equation}
This matches the serial formula \eqref{eq:residual} exactly.

\subsubsection{Thread safety}

\texttt{MPI\_Init\_thread} is called with the
\texttt{MPI\_THREAD\_FUNNELED} level, guaranteeing that only the master thread
issues MPI calls.  All \texttt{\#pragma omp parallel for collapse(2)} and
\texttt{reduction} pragmas are confined inside functions in
\texttt{jacobi\_omp\_mpi.h}; MPI calls in \texttt{main} remain strictly
outside any parallel region.

\subsubsection{Timing}

Each rank measures its own elapsed times with \texttt{MPI\_Wtime()}.
Before writing the CSV, \texttt{MPI\_Reduce(MPI\_MAX, root=0)} collects
the \emph{bottleneck} (slowest rank) time for each function.  
Only rank~0 opens and writes the output file, preventing concurrent file access.

% ============================================================
\section{Build System}
% ============================================================

All three targets are declared in a single \texttt{CMakeLists.txt}
with dependency sets that grow with the level of parallelism:

\begin{center}
\begin{tabular}{lll}
\toprule
CMake target       & Source file          & Linked libraries \\
\midrule
\texttt{jacobi}         & \texttt{jacobi.c}         & \texttt{m}, \texttt{rt} \\
\texttt{jacobi\_just\_omp}  & \texttt{jacobi\_just\_omp.c}  & \texttt{m}, \texttt{OpenMP::OpenMP\_C} \\
\texttt{jacobi\_omp\_mpi}   & \texttt{jacobi\_omp\_mpi.c}   & \texttt{m}, \texttt{OpenMP::OpenMP\_C}, \texttt{MPI::MPI\_C} \\
\bottomrule
\end{tabular}
\end{center}

\noindent
The global compiler flag \texttt{-O3} applies to all targets.
\texttt{find\_package(OpenMP REQUIRED)} and \texttt{find\_package(MPI REQUIRED)}
ensure portability across HPC toolchains without hard-coding paths.

% ============================================================
\section{Timing and Output Format}
% ============================================================

Every executable appends one row to a shared CSV file whose path is
read from the environment variable \texttt{TIMING\_CSV} (default:
\texttt{timing\_results.csv}).  The header is written only when the file
is empty, so repeated runs accumulate safely without duplication.

The CSV schema is:
\begin{verbatim}
case, nx, ny, number_omp_threads, number_mpi_processes,
fill_boundary_grid_time, fill_grid_without_boundary_time,
update_time, check_convergence_time, total_time, iterations
\end{verbatim}

\begin{itemize}
  \item \texttt{total\_time} spans from grid allocation to \texttt{free()}.
  \item \texttt{update\_time} and \texttt{check\_convergence\_time} are
        \emph{cumulative} across all iterations.
  \item For the MPI variant, all times are the per-function maximum over ranks
        (\texttt{MPI\_Reduce(MPI\_MAX)}), reflecting the true wall-clock
        bottleneck.
\end{itemize}

% ============================================================
\section{Benchmark Campaign via SLURM}
% ============================================================

The Python driver \texttt{main.py} automates the full benchmark pipeline:

\begin{enumerate}
  \item \textbf{CMake build} --- configures and compiles all targets in a
        \texttt{build/} subdirectory inside the scripts folder.
  \item \textbf{Parameter sweep} --- generates one SLURM batch script per
        combination of:
        \begin{itemize}
          \item grid sizes: $100\times100$ to $950\times950$ in steps of~50
                (19 sizes);
          \item OMP threads per rank ($T$): $1,2,4,8,16,32,64$;
          \item MPI ranks ($P$): $1,2,4,8,16,32,64$ (only for \texttt{omp\_mpi}).
        \end{itemize}
        The total job count is $19\cdot(1 + 7 + 7^2) = 19\cdot57 = 1083$ jobs.
  \item \textbf{Job monitoring} --- \texttt{sacct} is polled every 30\,s until
        every job reaches a terminal state.
  \item \textbf{Plot generation} --- calls \texttt{analysis.py} to produce all
        figures described in Section~\ref{sec:analysis}.
\end{enumerate}

All SLURM parameters (\texttt{SLURM\_PARTITION}, \texttt{SLURM\_ACCOUNT},
\texttt{SLURM\_TIME}) and the CSV path are overridable via environment
variables without modifying source code.  A \texttt{--dry-run} flag prints
all batch scripts without submitting, for inspection on a login node.

% ============================================================
\section{Performance Analysis}
\label{sec:analysis}
% ============================================================

\subsection{Methodology}

The Python module \texttt{analysis.py} reads the timing CSV with
\texttt{pandas} and produces one figure per timing column per grid size.
On every plot:
\begin{itemize}
  \item the \textbf{x-axis} is the number of OMP threads per MPI rank;
  \item each \textbf{line} corresponds to a unique \texttt{(case, \#MPI~ranks)}
        combination, averaged over repeated measurements at the same
        configuration;
  \item the \textbf{legend label} is formatted as \texttt{<case> (mpi=<P>)}.
\end{itemize}
This layout simultaneously exposes
\emph{strong intra-node scaling} (varying $T$ at fixed grid size) and
\emph{inter-node scaling} (comparing lines at different~$P$).

\subsection{Total time}

Figure~\ref{fig:total_100} through Figure~\ref{fig:total_950} show
\texttt{total\_time} for three representative grid sizes.
The serial baseline (label \texttt{none, mpi=1}) provides the reference.
Strong-scaling efficiency for the hybrid runs is defined as
\begin{equation}
  E(P,T) = \frac{t_{\text{serial}}}{P\cdot T\cdot t(P,T)},
\end{equation}
where $t_{\text{serial}}$ is the single-thread serial time and $t(P,T)$ is
the wall time with $P$~ranks and $T$~threads each.

\begin{figure}[htbp]
  \centering
  \includegraphics[width=0.82\textwidth]{../plots/total_time_100x100.png}
  \caption{Total wall-clock time vs.\ OMP thread count for a
           $100\times100$ grid.}
  \label{fig:total_100}
\end{figure}

\begin{figure}[htbp]
  \centering
  \includegraphics[width=0.82\textwidth]{../plots/total_time_500x500.png}
  \caption{Total wall-clock time vs.\ OMP thread count for a
           $500\times500$ grid.}
  \label{fig:total_500}
\end{figure}

\begin{figure}[htbp]
  \centering
  \includegraphics[width=0.82\textwidth]{../plots/total_time_950x950.png}
  \caption{Total wall-clock time vs.\ OMP thread count for a
           $950\times950$ grid.}
  \label{fig:total_950}
\end{figure}

\subsection{Kernel breakdown}

The four per-function timing plots allow identifying which phase limits scaling.

\subsubsection{\texttt{update} kernel}

The Jacobi sweep (equation~\eqref{eq:jacobi}) dominates the total runtime for
large grids because it is called at every iteration and performs $\mathcal{O}(N_x
N_y)$ floating-point operations with good spatial locality.
Figures~\ref{fig:update_100}--\ref{fig:update_950} show how its time varies
with parallelism.

\begin{figure}[htbp]
  \centering
  \begin{subfigure}[b]{0.48\textwidth}
    \includegraphics[width=\textwidth]{../plots/update_time_100x100.png}
    \caption{$100\times100$}
    \label{fig:update_100}
  \end{subfigure}
  \hfill
  \begin{subfigure}[b]{0.48\textwidth}
    \includegraphics[width=\textwidth]{../plots/update_time_950x950.png}
    \caption{$950\times950$}
    \label{fig:update_950}
  \end{subfigure}
  \caption{\texttt{update\_time}: cumulative time spent in the Jacobi sweep.
           For large grids the hybrid version is expected to scale well because
           each rank works on a reduced number of rows while OMP threads
           parallelise the inner loop.}
\end{figure}

\subsubsection{\texttt{check\_convergence} kernel}

The residual computation has the same $\mathcal{O}(N_x N_y)$ cost per call as
\texttt{update}, but the MPI variant introduces an
\texttt{MPI\_Allreduce} synchronisation barrier at every iteration.
For small grids this latency can dominate; for large grids the local
computation hides it.
Figures~\ref{fig:conv_100}--\ref{fig:conv_950} illustrate the crossover.

\begin{figure}[htbp]
  \centering
  \begin{subfigure}[b]{0.48\textwidth}
    \includegraphics[width=\textwidth]{../plots/check_convergence_time_100x100.png}
    \caption{$100\times100$}
    \label{fig:conv_100}
  \end{subfigure}
  \hfill
  \begin{subfigure}[b]{0.48\textwidth}
    \includegraphics[width=\textwidth]{../plots/check_convergence_time_950x950.png}
    \caption{$950\times950$}
    \label{fig:conv_950}
  \end{subfigure}
  \caption{\texttt{check\_convergence\_time}: cumulative time in the
           residual computation.  The MPI synchronisation overhead is visible
           as a floor that does not decrease with more OMP threads.}
\end{figure}

\subsubsection{Initialisation kernels}

\texttt{fill\_boundary\_grid} and \texttt{fill\_grid\_withouth\_boundary}
are called once and are therefore negligible fractions of the total time for
any non-trivial number of iterations.
They are included nevertheless to verify that the parallelisation of the
initialisation phase does not introduce overhead relative to the serial
baseline (Figures~\ref{fig:fill_bnd}--\ref{fig:fill_int}).

\begin{figure}[htbp]
  \centering
  \begin{subfigure}[b]{0.48\textwidth}
    \includegraphics[width=\textwidth]{../plots/fill_boundary_grid_time_500x500.png}
    \caption{\texttt{fill\_boundary\_grid}}
    \label{fig:fill_bnd}
  \end{subfigure}
  \hfill
  \begin{subfigure}[b]{0.48\textwidth}
    \includegraphics[width=\textwidth]{../plots/fill_grid_without_boundary_time_500x500.png}
    \caption{\texttt{fill\_grid\_without\_boundary}}
    \label{fig:fill_int}
  \end{subfigure}
  \caption{Initialisation times for a $500\times500$ grid.  Both are one-shot
           costs independent of the iteration count and are expected to be much
           smaller than the solver kernels.}
\end{figure}

\subsection{Expected scaling behaviour}

\paragraph{Serial vs.\ OpenMP.}
For $T$ threads Amdahl's law predicts a speed-up upper bound of
$S(T) \le 1/(f_s + (1-f_s)/T)$, where $f_s$ is the serial fraction
(boundary setup, \texttt{memcpy}, loop overhead).
For large grids $f_s \ll 1$ and near-linear scaling is expected up to
the point where the problem fits in the last-level cache.
Beyond that, memory bandwidth becomes the bottleneck and scaling saturates.

\paragraph{MPI strong scaling.}
With $P$ ranks each processing $N_x/P$ rows, the computation scales as
$\mathcal{O}(N_x N_y / P)$, while the halo-exchange cost per iteration is
$\mathcal{O}(N_y)$ (one row per neighbour).
For fixed $N_x N_y$, increasing $P$ reduces computation but keeps communication
constant, so an efficiency drop is expected as $P \to N_x$.

\paragraph{Hybrid MPI+OpenMP.}
The hybrid configuration is expected to outperform pure MPI for a fixed total
process count $P\cdot T$, because replacing $T$ MPI ranks within a node with
$T$ OMP threads eliminates intra-node message passing.
In practice the optimal split depends on memory-bandwidth saturation and
network latency; the plots in this section expose these trade-offs empirically.

% ============================================================
\section{Conclusions}
% ============================================================

The repository provides a self-contained benchmark that covers the full
lifecycle of an HPC study: from a cleanly formulated numerical problem
(2D Laplace / Jacobi iteration), through three implementations of increasing
parallelism (serial, shared-memory OMP, distributed+shared MPI+OMP), to an
automated benchmark campaign on a SLURM cluster and a quantitative
performance analysis.

The key findings anticipated from the experimental plots are:
\begin{enumerate}
  \item The \texttt{update} kernel dominates total time for large grids and
        is the primary beneficiary of both OMP and MPI parallelism.
  \item The \texttt{check\_convergence} residual under the MPI variant
        exhibits a communication floor set by \texttt{MPI\_Allreduce} latency,
        visible at small grid sizes.
  \item Hybrid configurations reduce intra-node traffic and are expected to
        achieve higher efficiency than pure MPI at equivalent total core counts
        on shared-memory nodes.
  \item Initialisation costs are negligible fractions of total time for the
        grid-size range explored.
\end{enumerate}

\end{document}
